\BeleskeID{08-s003}
% element: 08-s003:e04
% element: 08-s003:notes
Ulazni signal $a$ kolo tumači kao logičku jedinicu ili nulu, dok izlazni signal $b$ kolo generiše, postavlja ili daje kao jedinicu ili nulu. Logičko stanje $A$ određeno je tumačenjem ulaznog signala, a $B$ tumačenjem izlaznog signala. Zato logička funkcija $B=\overline A$ nije isto što i fizička funkcija $b=f(a)$, koja je po pravilu nelinearna.

Indeks $i$ dolazi od input i označava ulaz: fizičku veličinu koja nosi informaciju o nulama i jedinicama na ulazu. Indeks $o$ dolazi od output i označava istu ulogu na izlazu. U zapisima $a_i(t)$ i $b_o(t)$ eksplicitno se vidi i vremenska zavisnost.

Signal ne može trenutno da promeni svoje stanje. Mora proći kroz sve vrednosti između dva stanja čak i kada interval promene u pojednostavljenom modelu smatramo beskonačno kratkim. Gornja osenčena oblast pokazuje neizvesnost šta kolo prepoznaje tokom ulazne ivice; donja pokazuje šta i kada daje na izlazu.
